\section{Related Work}\label{sec:related_work}

\subsection{Large-Vocabulary and Structured Chord Recognition}\label{subsec:rw_large_vocab}

Early ACR systems restricted classification to 24 major/minor triads using chroma vectors and hidden Markov models~\cite{korzeniowski2016feature,pauwels2019twenty}. To scale to larger vocabularies without combinatorial explosion, McFee and Bello~\cite{mcfee2017structured} formulated a structured multi-task loss over 170 classes (14 qualities), regularizing predictions via root, bass, and pitch-class targets; however, their taxonomy maps all upper extensions ($9\text{th}, 11\text{th}, 13\text{th}$) onto basic triads or tetrads.

Jiang et al.~\cite{jiang2019large} preserved upper tertian degrees by decomposing chords into six categorical heads (Root+Triad, Bass, $7\text{th}, 9\text{th}, 11\text{th}, 13\text{th}$) trained with CRNNs, inverse-frequency class reweighting, and CRF decoding. Self-attention architectures further improved temporal context modeling: the Bidirectional Transformer for Chords (BTC)~\cite{park2019btc} replaced recurrent layers with bidirectional self-attention, and Kim and Park~\cite{kim2026event} introduced event-level boundary pooling to suppress frame oversegmentation on non-triad harmonies.

Akram et al.~\cite{akram2025chordformer} proposed ChordFormer, pairing Macaron-style Conformer blocks~\cite{gulati2020conformer} with the 6-head decomposition of Jiang et al.~\cite{jiang2019large} to capture local harmonic partials via depthwise convolution and global phrase structure via multi-head attention. Li et al.~\cite{li2025large} applied contrastive learning and Noisy Student pseudo-labeling to unannotated audio across 170 classes. However, teacher-student pseudo-labeling suffers from confirmation bias on the harmonic tail: because teacher networks exhibit near-zero recall on rare alterations, pseudo-labels reinforce majority diatonic priors. Our approach instead supplies noise-free, symbolically verified ground-truth audio across the full extension vocabulary.

\subsection{Synthetic Audio Datasets in MIR}\label{subsec:rw_synthetic_audio}

Synthetic multitrack rendering offers a scalable remedy for annotation scarcity in MIR. In source separation, Slakh2100~\cite{manilow2019cutting} rendered $2{,}100$ MIDI arrangements through virtual instrument libraries to provide isolated stems. For chord recognition, Ostermann et al.~\cite{ostermann2023aam} developed Artificial Audio Multitracks (AAM), comprising $3{,}000$ FluidSynth-rendered tracks generated via algorithmic composition in JFugue.

Majchrzak and Ma{\'n}dziuk~\cite{majchrzak2025training} benchmarked synthetic pretraining on BTC and Harmony Transformer models using AAM, showing consistent accuracy gains on major/minor pop progressions after real-audio finetuning, but zero transfer to extended chord vocabularies and degraded accuracy on jazz repertoires. Because AAM samples progressions exclusively from unextended diatonic triads ($\text{I}$--$\text{vi}$)~\cite{ostermann2023aam}, it replicates the triadic skew of real corpora. In contrast, our generator acts as an explicit distributional intervention: upweighting rare triads and altered extensions via quota scheduling, enforcing physical voicings so labeled extensions sound acoustically, and applying random-mix augmentation to bridge the SoundFont-to-studio gap.
